%%%%%%%%%%%%%%%%%%%%%%%%%%%%%%%%%%%%%%%%%%%%%%%%%%%%%%%%%%%%%%%%%%%%%%%%%%%%%%%%%%%%%%%%
% Projeto de Extensão da disciplina de Programação 3
% Sub-projeto do ProgramAuto
% Autores:
%     Prof. Dr. Ruben Carlo Benante
%     Nome do Aluno 1
%     Nome do Aluno 2
%
% Data: 2026-09-20
%
% Assunto: Bibliotecas graficas, IHM, GTK, Ogre, QT, QMake e Make
%
%%%%%%%%%%%%%%%%%%%%%%%%%%%%%%%%%%%%%%%%%%%%%%%%%%%%%%%%%%%%%%%%%%%%%%%%%%%%%%%%%%%%%%%%


%%%%%%%%%%%%%%%%%%%%%%%%%%%%%%%%%%%%%%%%%%%%%%%%%%%%%%%%%%%%%%%%%%%%%%%%%%%%%%%%%%%%%%%%
% Para gerar o PDF use uma das 2 opções abaixo:
%
% Opção 1: com makefile, comando único.
%    $ make ext-programaX-benante-sobrenome1-sobrenome2.pdf
%
% Opção 2: com os 3 comandos diretamente:
%    $ pdflatex pext-programaX-benante-sobrenome1-sobrenome2.tex -o pext-programaX-benante-sobrenome1-sobrenome2.pdf
%    $ bibtex biblio
%    $ pdflatex pext-programaX-benante-sobrenome1-sobrenome2.tex -o pext-programaX-benante-sobrenome1-sobrenome2.pdf


%%%%%%%%%%%%%%%%%%%%%%%%%%%%%%%%%%%%%%%%%%%%%%%%%%%%%%%%%%%%%%%%%%%%%%%%%%%%%%%%%%%%%%%%
% preambulo %%%%%%%%%%%%%%%%%%%%%%%%%%%%%%%%%%%%%%%%%%%%%%%%%%%%%%%%%%%%%%%%%%%%%%%%%%%%
\documentclass[a4paper,12pt]{article} %twocolumn
\usepackage[left=2.5cm,right=2cm,top=2.5cm,bottom=2cm]{geometry}
\usepackage[utf8]{inputenc} % letras acentuadas
\usepackage[portuguese]{babel} % tradução de títulos
\usepackage[colorlinks]{hyperref}
\usepackage{algorithm} % ambiente para índice de algoritmos
\usepackage{algpseudocode} % fonte e estilo do algoritmo
\usepackage{graphicx}
\usepackage{indentfirst}
\usepackage{url}
% \usepackage{natbib}
%[noend]

\DeclareUrlCommand\email{\urlstyle{mm}} % comando para email bonito
\floatname{algorithm}{Algoritmo} % tradução da palavra algorítimo no ambiente de índice


%%%%%%%%%%%%%%%%%%%%%%%%%%%%%%%%%%%%%%%%%%%%%%%%%%%%%%%%%%%%%%%%%%%%%%%%%%%%%%%%%%%%%%%%
% capa %%%%%%%%%%%%%%%%%%%%%%%%%%%%%%%%%%%%%%%%%%%%%%%%%%%%%%%%%%%%%%%%%%%%%%%%%%%%%%%%%
\title{Programa 3: Interface Pessoa-Maquina com C++}
\author{Ruben Carlo Benante \\ Autor2 \\ Autor3}

\begin{document}

\maketitle


%%%%%%%%%%%%%%%%%%%%%%%%%%%%%%%%%%%%%%%%%%%%%%%%%%%%%%%%%%%%%%%%%%%%%%%%%%%%%%%%%%%%%%%%
% resumo %%%%%%%%%%%%%%%%%%%%%%%%%%%%%%%%%%%%%%%%%%%%%%%%%%%%%%%%%%%%%%%%%%%%%%%%%%%%%%%

\begin{abstract}

\textbf{Assunto:} Ensino da Linguagem de Programação \texttt{C} / \texttt{CPP}.

% descrever em poucas palavras seu projeto aqui
Este projeto de extensao aborda o ensino e desenvolvimento de interfaces pessoa-maquina abordando bibliotecas graficas e compilacao de softwares em C++. Por meio de seis videos educativos, visa-se capacitar a comunidade academica na construcao de softwares visuais.

\textbf{Local:} Escola Politécnica de Pernambuco - UPE/POLI

\textbf{Órgão Financiador:} N/A

\textbf{Caracterização:} Projeto de Extensão requisito da disciplina de DCExt Programação 3, sub-projeto integrante do Projeto \texttt{ProgramAuto}

% Este é o fim do resumo.

\end{abstract}


%%%%%%%%%%%%%%%%%%%%%%%%%%%%%%%%%%%%%%%%%%%%%%%%%%%%%%%%%%%%%%%%%%%%%%%%%%%%%%%%%%%%%%%%
% artigo %%%%%%%%%%%%%%%%%%%%%%%%%%%%%%%%%%%%%%%%%%%%%%%%%%%%%%%%%%%%%%%%%%%%%%%%%%%%%%%
% seção de introdução %%%%%%%%%%%%%%%%%%%%%%%%%%%%%%%%%%%%%%%%%%%%%%%%%%%%%%%%%%%%%%%%%%
\section{Introdução}

% Descrever melhor seu projeto aqui

O ensino da Linguagem de Programação \texttt{C} / \texttt{CPP} expandiu-se da manipulacao em linha de comando para a construcao de sistemas visuais interativos e graficos de alto desempenho.

Diante do desafio de projetar interfaces homem-maquina (IHM) responsivas e acessiveis, torna-se essencial orientar os estudantes na escolha das ferramentas adequadas, diferenciando bibliotecas focadas em multimida/jogos de frameworks voltados a aplicacoes industriais.\ldots

O ambiente em C++ abrange desde o controle de baixo nivel de janelas ate a automacao do fluxo de compilacao com QMake e Make. \ldots

%%%%%%%%%%%%%%%%%%%%%%%%%%%%%%%%%%%%%%%%%%%%%%%%%%%%%%%%%%%%%%%%%%%%%%%%%%%%%%%%%%%%%%%%
% seção de objetivos %%%%%%%%%%%%%%%%%%%%%%%%%%%%%%%%%%%%%%%%%%%%%%%%%%%%%%%%%%%%%%%%%%%
\section{Objetivos}

\subsection{Objetivo Geral}

%Descrever o objetivo geral a ser alcançado

Criação e divulgacao de seis videos instrucionais focados no ensino de bibliotecas graficas em C++, cobrindo o ciclo completo desde os conceitos ate a automacao de build de aplicacoes Qt com arquivos \texttt{.ui} e aruivos de projeto \texttt{.pro}.

\subsection{Objetivos Específicos}

% Listar os objetivos específicos

Como objetivos específicos este projeto visa atingir os seguintes pontos, conforme abaixo:

\begin{itemize}
    \item Apresentar os conceitos fundamentais de IHM, usabilidade e a comunicacao com a GPU via OpenGL/Vulkan/DirectX;
    \item Demonstrar a estrutura de loop principal do jogo e manipulacao 2D com a biblioteca Allegro 5;
    \item Comparar e integrar o GTK e o Ogre para renderizacao;
    \item Explicar a estrutura do QT Framework e o sistema de Signals e Slots;
    \item Ensinar o uso do Qt Designer, o que sao os arquivos \texttt{.ui} e a integracao com o codigo C++ utilizando a ferramenta \texttt{uic};
    \item Automatizar a compilacao de softwares em C++ utilizando arquivos de projeto \texttt{.pro}, QMake e Make.
\end{itemize}


%%%%%%%%%%%%%%%%%%%%%%%%%%%%%%%%%%%%%%%%%%%%%%%%%%%%%%%%%%%%%%%%%%%%%%%%%%%%%%%%%%%%%%%%
% seção de justificativa %%%%%%%%%%%%%%%%%%%%%%%%%%%%%%%%%%%%%%%%%%%%%%%%%%%%%%%%%%%%%%%
\section{Justificativa}

%Justificar aqui o seu projeto.

As características tal e tal da linguagem C (ou CPP se aluno de Programação 3) são de fundamental importância para se atingir a excelência e tal e coisa e tal, conforme denota \cite{Benante2008phd}.

A Figura \ref{fig:graficobarra} abaixo identifica a principal razão, exposta acima, para bla bla

\begin{figure}[ht]
\centering
\includegraphics[width=.60\linewidth]{imagens/exemplo1.png}
\caption{Exemplo de gráfico em barras}
\label{fig:graficobarra}
\end{figure}


%%%%%%%%%%%%%%%%%%%%%%%%%%%%%%%%%%%%%%%%%%%%%%%%%%%%%%%%%%%%%%%%%%%%%%%%%%%%%%%%%%%%%%%%
% seção de metodologia %%%%%%%%%%%%%%%%%%%%%%%%%%%%%%%%%%%%%%%%%%%%%%%%%%%%%%%%%%%%%%%%%
\section{Metodologia}

% Descrever como (por quais métodos) os objetivos serão alcançados.

Este projeto contará com a criação, roteiro, produção, gravação e divugação de quatro (4) vídeos tal e tal.

Os videos serão divididos nos conteúdos tal e tal.

A abordagem dada é tal e tal e coisa e tal.

%%%%%%%%%%%%%%%%%%%%%%%%%%%%%%%%%%%%%%%%%%%%%%%%%%%%%%%%%%%%%%%%%%%%%%%%%%%%%%%%%%%%%%%%
% subseção equipamentos %%%%%%%%%%%%%%%%%%%%%%%%%%%%%%%%%%%%%%%%%%%%%%%%%%%%%%%%%%%%%%%%
\subsection{Equipamentos Necessários}

%Listar os equipamentos necessários para a implementação do projeto.

Para realizar este projeto é preciso tal e tal, além de tal, e acesso à internet, e tal e tal.

Além desses equipamentos, também é preciso coisa e tal.


%%%%%%%%%%%%%%%%%%%%%%%%%%%%%%%%%%%%%%%%%%%%%%%%%%%%%%%%%%%%%%%%%%%%%%%%%%%%%%%%%%%%%%%%
% subseção com algoritmo %%%%%%%%%%%%%%%%%%%%%%%%%%%%%%%%%%%%%%%%%%%%%%%%%%%%%%%%%%%%%%%
\subsection{Implementação}

A implementação do projeto se dará de forma online, em uma \textit{playlist} numa plataforma de reprodução de videos da escolha do professor orientador, e a mesma será divulgada tanto para alunos da Escola Politécnica de Pernambuco (POLI/UPE) quanto para a comunidade em geral através dos canais de comunicação disponíveis.



%%%%%%%%%%%%%%%%%%%%%%%%%%%%%%%%%%%%%%%%%%%%%%%%%%%%%%%%%%%%%%%%%%%%%%%%%%%%%%%%%%%%%%%%
% seção Plano de Trabalho %%%%%%%%%%%%%%%%%%%%%%%%%%%%%%%%%%%%%%%%%%%%%%%%%%%%%%%%%%%%%%
\section{Plano de Trabalho}

% Esta seção estabelece as atividades a serem realizadas, organizadas em etapas, para depois serem alocadas na tabela na seção seguinte, cronograma.

% dicas de etapas: desenvolvimento de roteiros, videos, edição, gravação.

Este projeto está subdividido em XXX etapas, que contemplam ... bla bla
e estão organizadas em ordem cronológica na Tabela \ref{tab:cronograma}.

São elas:

\begin{description}
 \item[Etapa 1: Projeto.] Organização do projeto, composição das idéias, seleção de tema, discussão com membros do grupo, divisão de tarefas e atividades administrativas relevantes.
 \item[Etapa 2: Pesquisa.] (descreva a etapa)
 \item[Etapa 3: Conteúdo.] (descreva a etapa)
 \item[Etapa 4: Desenvolvimento.] (descreva a etapa)
 \item[Etapa 5: Revisão.] (descreva a etapa)
 \item[Etapa 6: Validação.] (descreva a etapa)
\end{description}


%%%%%%%%%%%%%%%%%%%%%%%%%%%%%%%%%%%%%%%%%%%%%%%%%%%%%%%%%%%%%%%%%%%%%%%%%%%%%%%%%%%%%%%%
% seção Plano de Trabalho %%%%%%%%%%%%%%%%%%%%%%%%%%%%%%%%%%%%%%%%%%%%%%%%%%%%%%%%%%%%%%
\section{Cronograma}

%Em conjunto com a seção de Plano de Trabalho, a seção de cronograma coloca as atividades dispostas numa linha do tempo. Utilize uma tabela para melhor visualização.

Segue abaixo o cronograma com tal e tal

\begin{table}[H]
\begin{center}
 \caption{Tabela de Cronograma do Projeto de Extensão}
\begin{tabular}{|l|c|c|c|c|c|c|}
  \hline \hline
   \textbf{Etapas} & \textbf{Sem. 1} & \textbf{Sem. 2} & \textbf{Sem. 3} & \textbf{Sem. 4} & \textbf{Sem. 5} & \textbf{Sem. 6} \\ \hline \hline
   1. Projeto. & X & X &  &  &  &  \\ \hline
   2. Pesquisa. &  &  &  &  &  &  \\ \hline
   3. Conteúdo. &  &  &  &  &  &  \\ \hline
   4. Desenvolvimento. &  &  &  &  &  &  \\ \hline
   5. Revisão. &  &  &  &  &  &  \\ \hline
   6. Validação. &  &  &  &  &  &  \\ \hline \hline
\end{tabular}
\label{tab:cronograma}
\end{center}
\end{table}


%%%%%%%%%%%%%%%%%%%%%%%%%%%%%%%%%%%%%%%%%%%%%%%%%%%%%%%%%%%%%%%%%%%%%%%%%%%%%%%%%%%%%%%%
% seção de impactos alcançados %%%%%%%%%%%%%%%%%%%%%%%%%%%%%%%%%%%%%%%%%%%%%%%%%%%%%%%%%
\section{Impactos e Transferências}


%%%%%%%%%%%%%%%%%%%%%%%%%%%%%%%%%%%%%%%%%%%%%%%%%%%%%%%%%%%%%%%%%%%%%%%%%%%%%%%%%%%%%%%%
% subseção de impacto científico %%%%%%%%%%%%%%%%%%%%%%%%%%%%%%%%%%%%%%%%%%%%%%%%%%%%%%%
\subsection{Impacto Científico}

Não há impacto científico relevante (N/A).


%%%%%%%%%%%%%%%%%%%%%%%%%%%%%%%%%%%%%%%%%%%%%%%%%%%%%%%%%%%%%%%%%%%%%%%%%%%%%%%%%%%%%%%%
% subseção de impacto tecnológico %%%%%%%%%%%%%%%%%%%%%%%%%%%%%%%%%%%%%%%%%%%%%%%%%%%%%%
\subsection{Impacto Tecnológico}

Não há impacto tecnológico relevante (N/A).


%%%%%%%%%%%%%%%%%%%%%%%%%%%%%%%%%%%%%%%%%%%%%%%%%%%%%%%%%%%%%%%%%%%%%%%%%%%%%%%%%%%%%%%%
% subseção de econônimo %%%%%%%%%%%%%%%%%%%%%%%%%%%%%%%%%%%%%%%%%%%%%%%%%%%%%%%%%%%%%%%%
\subsection{Impacto Econômico}

Não há impacto econômico relevante (N/A).


%%%%%%%%%%%%%%%%%%%%%%%%%%%%%%%%%%%%%%%%%%%%%%%%%%%%%%%%%%%%%%%%%%%%%%%%%%%%%%%%%%%%%%%%
% subseção de impacto social %%%%%%%%%%%%%%%%%%%%%%%%%%%%%%%%%%%%%%%%%%%%%%%%%%%%%%%%%%%
\subsection{Impacto Social}

O projeto visa contribuir com a sociedade na forma de... bla ... bla... bla


%%%%%%%%%%%%%%%%%%%%%%%%%%%%%%%%%%%%%%%%%%%%%%%%%%%%%%%%%%%%%%%%%%%%%%%%%%%%%%%%%%%%%%%%
% subseção de impacto ambiental %%%%%%%%%%%%%%%%%%%%%%%%%%%%%%%%%%%%%%%%%%%%%%%%%%%%%%%%
\subsection{Impacto Ambiental}

Não há impacto ambiental relevante (N/A).


%%%%%%%%%%%%%%%%%%%%%%%%%%%%%%%%%%%%%%%%%%%%%%%%%%%%%%%%%%%%%%%%%%%%%%%%%%%%%%%%%%%%%%%%
% subseção de transferências %%%%%%%%%%%%%%%%%%%%%%%%%%%%%%%%%%%%%%%%%%%%%%%%%%%%%%%%%%%
\subsection{Transferências}

O ProgramAuto terá como objetivo a execução de aulas para ensinar a linguagem C aos discentes interessados. Contará com sua exposição de ensino gravada que será disponibilizada e a elaboração de relatórios a fim de cumprir com os aspectos estabelecidos no Plano de Trabalho.

Este projeto, parte integrante do ProgramAuto, visa ...

%O que mais o seu projeto agrega? O que é transferido? De onde vem? Para onde vai?


%%%%%%%%%%%%%%%%%%%%%%%%%%%%%%%%%%%%%%%%%%%%%%%%%%%%%%%%%%%%%%%%%%%%%%%%%%%%%%%%%%%%%%%%
% seção de resultados esperados %%%%%%%%%%%%%%%%%%%%%%%%%%%%%%%%%%%%%%%%%%%%%%%%%%%%%%%%
\section{Resultados Esperados}

É esperado que após a execução com sucesso deste projeto,

%%%%%%%%%%%%%%%%%%%%%%%%%%%%%%%%%%%%%%%%%%%%%%%%%%%%%%%%%%%%%%%%%%%%%%%%%%%%%%%%%%%%%%%%
% Autores %%%%%%%%%%%%%%%%%%%%%%%%%%%%%%%%%%%%%%%%%%%%%%%%%%%%%%%%%%%%%%%%%%%%%%%%%%%%%%
\section*{Detalhamento dos Autores}

%%%%%%%%%%%%%%%%%%%%%%%%%%%%%%%%%%%%%%%%%%%%%%%%%%%%%%%%%%%%%%%%%%%%%%%%%%%%%%%%%%%%%%%%
% Discentes %%%%%%%%%%%%%%%%%%%%%%%%%%%%%%%%%%%%%%%%%%%%%%%%%%%%%%%%%%%%%%%%%%%%%%%%%%%%
\subsection*{Discentes}

\begin{enumerate}
    \item \textbf{Nome Completo:} Fulano de tal
    \begin{description}
        \item [Email:] \email{blabla@poli.br}
        \item [Endereço:]
        \item [Matrícula:]
        \item [CPF:]
        \item [RG:]
        \item [Telefone:]
        \item [Currículo Lattes:] \url{http://lattes.cnpq.br/nnnnn}
    \end{description}

    \item \textbf{Nome Completo:} Fulano de tal
    \begin{description}
        \item [Email:] \email{blabla@poli.br}
        \item [Endereço:]
        \item [Matrícula:]
        \item [CPF:]
        \item [RG:]
        \item [Telefone:]
        \item [Currículo Lattes:] \url{http://lattes.cnpq.br/nnnnn}
    \end{description}

    \item \textbf{Nome Completo:} Fulano de tal
    \begin{description}
        \item [Email:] \email{blabla@poli.br}
        \item [Endereço:]
        \item [Matrícula:]
        \item [CPF:]
        \item [RG:]
        \item [Telefone:]
        \item [Currículo Lattes:] \url{http://lattes.cnpq.br/nnnnn}
    \end{description}

    \item \textbf{Nome Completo:} Fulano de tal
    \begin{description}
        \item [Email:] \email{blabla@poli.br}
        \item [Endereço:]
        \item [Matrícula:]
        \item [CPF:]
        \item [RG:]
        \item [Telefone:]
        \item [Currículo Lattes:] \url{http://lattes.cnpq.br/nnnnn}
    \end{description}

%     \item \textbf{Nome Completo:} Fulano de tal
%     \begin{description}
%         \item [Email:] \email{blabla@poli.br}
%         \item [Endereço:]
%         \item [Matrícula:]
%         \item [CPF:]
%         \item [RG:]
%         \item [Telefone:]
%         \item [Currículo Lattes:] \url{http://lattes.cnpq.br/nnnnn}
%     \end{description}
\end{enumerate}


%%%%%%%%%%%%%%%%%%%%%%%%%%%%%%%%%%%%%%%%%%%%%%%%%%%%%%%%%%%%%%%%%%%%%%%%%%%%%%%%%%%%%%%%
% Docentes %%%%%%%%%%%%%%%%%%%%%%%%%%%%%%%%%%%%%%%%%%%%%%%%%%%%%%%%%%%%%%%%%%%%%%%%%%%%%
\subsection*{Docentes}

\begin{enumerate}
    \item \textbf{Nome Completo:} Ruben Carlo Benante
    \begin{description}
        \item [Email:] \email{rcb@poli.br}
        \item [Matrícula:] 11238-0
        \item [Currículo Lattes:] \url{http://lattes.cnpq.br/3366717378277623}
    \end{description}

%     \item \textbf{Nome Completo:}
%     \begin{description}
%         \item [Email:] \email{no@no.no}
%         \item [Matrícula:]
%         \item [Currículo Lattes:]
%     \end{description}
\end{enumerate}


%%%%%%%%%%%%%%%%%%%%%%%%%%%%%%%%%%%%%%%%%%%%%%%%%%%%%%%%%%%%%%%%%%%%%%%%%%%%%%%%%%%%%%%%
% referências bibliográficas %%%%%%%%%%%%%%%%%%%%%%%%%%%%%%%%%%%%%%%%%%%%%%%%%%%%%%%%%%%
%\section*{Referências Bibliográficas}

% cite todos, mesmo os não referenciados %%%%%%%%%%%%%%%%%%%%%%%%%%%%%%%%%%%%%%%%%%%%%%%
\nocite{*}


%%%%%%%%%%%%%%%%%%%%%%%%%%%%%%%%%%%%%%%%%%%%%%%%%%%%%%%%%%%%%%%%%%%%%%%%%%%%%%%%%%%%%%%%
% se necessario %%%%%%%%%%%%%%%%%%%%%%%%%%%%%%%%%%%%%%%%%%%%%%%%%%%%%%
% troca autor and autor por autor & autor, na bibliografia. O dcu usa "and"
%\renewcommand{\harvardand}{\&} % troca and pro &. O dcu usa "and"

% Estilos de bibliografia %%%%%%%%%%%%%%%%%%%%%%%%%%%%%%%%%%%%%%%%%%%%%%%%%%%%%%%%%%%%%%
% \bibliographystyle{abnt-alf} % Estilo alfabético da ABNT. Opção [num] para estilo numérico
% \bibliographystyle{apalike}
% \bibliographystyle{dcu} %citacao como (autor and autor, ano). Parece apalike. Rev. Control. Automacao. Use com harvard
% \bibliographystyle{agsm} % padrao harvard fica (autor & autor ano).
\bibliographystyle{acm}

%%%%%%%%%%%%%%%%%%%%%%%%%%%%%%%%%%%%%%%%%%%%%%%%%%%%%%%%%%%%%%%%%%%%%%%%%%%%%%%%%%%%%%%%
% arquivo de banco de dados das referências %%%%%%%%%%%%%%%%%%%%%%%%%%%%%%%%%%%%%%%%%%%%
% renomear para o número do exercício correto
% o arquivo de bibliografia pode se chamar qualquer coisa, isso não muda o comando de gerar o PDF.
% Por exemplo para 'mybiblio.bib', use \bibliography{mybiblio} e os comandos pdflatex e bibtex continuam os mesmos identicos com exN.
\bibliography{biblio}

\end{document}
